\documentclass[11pt]{article}
\usepackage[margin=1in]{geometry}
\usepackage{amsmath,amssymb}
\usepackage[hidelinks]{hyperref}

\newcommand{\R}{\mathbb R}
\newcommand{\D}{D}

\title{Literature-implied answer for a wedge Laplacian domain question}
\author{Run \texttt{fa\_banach\_001}, lane 17}
\date{2026-07-18}

\begin{document}
\maketitle

\section*{Source Signal}

The source paper is Emiel Lorist and Mark Veraar,
\emph{Singular stochastic integral operators}, arXiv:1902.10620.
In the appendix, after Proposition A.2 on the heat semigroup on a wedge
\[
\D=\{(r\cos\varphi,r\sin\varphi): r>0,\ 0<\varphi<\kappa\},
\]
the authors consider the Dirichlet Laplacian on
\[
X=L^q(\D, |x|^{\theta-2}).
\]
They prove sectoriality for
\[
  -\frac{\pi}{\kappa}<\frac{\theta}{q}<2+\frac{\pi}{\kappa},
\]
but state that the domain is not known on the full range.  They quote the
known formula
\[
D(\Delta)=\{u: u,\ u/|x|^2,\ \partial^\alpha u\in L^q(\D,|x|^{\theta-2})
\text{ for }|\alpha|=2\}
\]
only for
\[
  2-\frac{\pi}{\kappa}<\frac{\theta}{q}<2+\frac{\pi}{\kappa}.
\]

The same source also contains a separate signal about whether condition
$(C_p)$ is \(p\)-independent.  That is a same-paper false positive: the source
paper says this was open and proves the independence immediately below.

\section*{Supporting Literature}

The supporting paper is Petru A. Cioica-Licht, Emiel Lorist, and P. Tobias
Werner, \emph{\(H^\infty\)-calculus for the Dirichlet Laplacian on conical
domains}, arXiv:2509.09519v2.  Its abstract states that it proves bounded
\(H^\infty\)-calculus for the Dirichlet Laplacian on conical domains and
wedges in weighted \(L^p\)-spaces with weights based on both the boundary
distance and the distance to the tip/edge.

Theorem A, equivalently Theorem 4.1 in the body of the paper, defines the
Dirichlet Laplacian on
\[
  L^p(\mathcal D,w_{\gamma,\nu}^{\mathcal D}),\qquad
  w_{\gamma,\nu}^{\mathcal D}=\rho_\circ^\nu
  \left(\frac{\rho_{\mathcal D}}{\rho_\circ}\right)^\gamma,
\]
with domain \(\dot W^{2,p}_{\rm Dir}(\mathcal D,w_{\gamma,\nu}^{\mathcal D})
\cap L^p(\mathcal D,w_{\gamma,\nu}^{\mathcal D})\).  It proves closedness,
graph-norm equivalence, a homogeneous domain identification, and, in the
natural parameter range, bounded \(H^\infty\)-calculus of angle zero.

\section*{Parameter Match}

For a planar wedge, \(d=2\) and the spherical cross-section is the interval
\((0,\kappa)\).  The eigenvalues of the one-dimensional Dirichlet
Laplace--Beltrami operator are \((n\pi/\kappa)^2\), so
\[
  \lambda_n^*=\frac{n\pi}{\kappa}.
\]
To recover the source paper's space \(L^q(\D,|x|^{\theta-2})\), take
\[
  p=q,\qquad \gamma=0,\qquad \nu=\theta-2.
\]
Then \(w_{\gamma,\nu}^{\D}=|x|^{\theta-2}\).  The \(H^\infty\)-calculus
range in the supporting paper becomes
\[
  -\frac{\pi}{\kappa}<\frac{\theta}{q}<2+\frac{\pi}{\kappa},
\]
which is exactly the sectoriality range quoted in arXiv:1902.10620.
The supporting theorem excludes the resonant values
\[
  \frac{\theta}{q}=2\pm \frac{n\pi}{\kappa},\qquad n\in\mathbb N.
\]

\section*{Conclusion}

The later Cioica-Licht--Lorist--Werner theorem gives the requested domain
description for the source paper's wedge Dirichlet Laplacian problem across
the sectorial range, except at the explicit resonant parameters above.  Thus
the target is not a new open problem for this run; it is a literature-implied
answer with partial scope.

\section*{References}

\begin{thebibliography}{9}

\bibitem{LV}
E. Lorist and M. Veraar,
\emph{Singular stochastic integral operators},
Analysis \& PDE 16 (2023), no. 2, 387--466;
arXiv:1902.10620.
\url{https://arxiv.org/abs/1902.10620}

\bibitem{CLLW}
P. A. Cioica-Licht, E. Lorist, and P. T. Werner,
\emph{\(H^\infty\)-calculus for the Dirichlet Laplacian on conical domains},
arXiv:2509.09519v2, 2025.
\url{https://arxiv.org/abs/2509.09519}

\end{thebibliography}

\end{document}

